% Generated draft: requires mathematical and source review.
\documentclass[11pt]{article}
\usepackage[margin=25mm]{geometry}
\usepackage{fontspec}
\setmainfont{Latin Modern Roman}
\usepackage{amsmath,amssymb,mathtools,unicode-math}
\setmathfont{Latin Modern Math}
\usepackage{enumitem,needspace,xurl,hyperref}
\setlength{\parindent}{0pt}
\setlength{\parskip}{5pt}
\setlength{\emergencystretch}{3em}
\newcommand{\N}{\mathbb{N}}
\newcommand{\Z}{\mathbb{Z}}
\newcommand{\Q}{\mathbb{Q}}
\newcommand{\R}{\mathbb{R}}
\newcommand{\C}{\mathbb{C}}
\newcommand{\F}{\mathbb{F}}
\newcommand{\A}{\mathbb{A}}
\newcommand{\PP}{\mathbb{P}}
\newcommand{\E}{\mathbb{E}}
\newcommand{\eps}{\varepsilon}
\newcommand{\dd}{\,\mathrm{d}}
\begin{document}
\section*{UnsolvedMath ID: 30002831}
\textbf{Source problem code:} catalogue entry
\subsection{Definitions and mathematical statement}
The principal target is the bijection in the supplied clean statement. The ground field is not specified in the record; the formulation below is conditional on the conventional interpretation over \(\C\). The precise conventions for deformation families and combinatorial mutation must be checked against the cited sources.

Let \(N\cong\Z^2\) and \(N_{\R}=N\otimes_{\Z}\R\). A Fano polygon is a two-dimensional convex lattice polygon \(P\subset N_{\R}\) whose vertices are primitive lattice vectors and whose strict interior contains the origin. Let \(X_P\) be the toric surface whose fan consists of the cones over the faces of \(P\). Write \(P\sim_{\mathrm{mut}}Q\) when \(P\) and \(Q\) are related by a finite sequence of combinatorial mutations and lattice isomorphisms, using the mutation convention of the cited conjectures.

A del Pezzo surface here means a normal projective surface \(X\) with \(K_X\) \(\Q\)-Cartier and \(-K_X\) ample. Its singularities are required to be cyclic quotient singularities. A \(\Q\)-Gorenstein deformation, abbreviated qG-deformation, is a flat deformation satisfying the relative \(\Q\)-Gorenstein condition; the exact admissibility and base-change conditions on such families are not supplied. A surface is locally qG-rigid when its singularity germs admit no nontrivial qG-deformations. Under the usual meaning of class \(\mathrm{TG}\), a surface belongs to this class if it admits a qG-degeneration to a toric del Pezzo surface. This meaning, including restrictions on the degeneration, requires source confirmation. Write \(X\sim_{\mathrm{qG}}Y\) for qG-deformation equivalence, with the permitted families and intermediate fibers understood according to the source convention.

Let \(\mathcal F\) be the collection of Fano polygons, allowing lattice isomorphisms, and let \(\mathcal D\) be the collection of locally qG-rigid class-\(\mathrm{TG}\) del Pezzo surfaces with cyclic quotient singularities. The principal conjecture asks whether
\[
\mathcal F/\!\sim_{\mathrm{mut}}
\;\longrightarrow\;
\mathcal D/\!\sim_{\mathrm{qG}},
\qquad
[P]\longmapsto[X],
\]
is a well-defined bijection, where \(X\) is a generic qG-deformation of \(X_P\). In particular, this assertion requires such a generic deformation to belong to \(\mathcal D\), its qG-deformation class to be independent of the generic choice, every class in \(\mathcal D/\!\sim_{\mathrm{qG}}\) to occur, and, for all Fano polygons \(P,Q\) and corresponding generic deformations \(X,Y\),
\[
P\sim_{\mathrm{mut}}Q
\quad\Longleftrightarrow\quad
X\sim_{\mathrm{qG}}Y.
\]
The record does not specify the deformation space or generic locus intended by ``generic'', so these requirements are part of the conjectural formulation rather than established existence or uniqueness claims.

The supplied original statement also records three companion assertions, which must be retained as context:
\begin{enumerate}
\item For every Fano polygon \(P\) and a generic qG-deformation \(X\) of \(X_P\), the maximally-mutable Laurent-polynomial space \(L_P^T\) is affine-linearly isomorphic, by a mirror map, to the space \(H_X^{\mathrm{ts}}\) of twisted-sector parameters of age strictly less than \(1\). The mirror map is asserted to identify the regularized quantum period with the classical period. The record does not define the superscript \(T\), the coefficient restrictions on \(L_P^T\), the precise twisted-sector convention, or the period normalizations.
\item Equal periods for Fano polygons with the same singularity content are asserted to imply mutation equivalence. A polygon alone does not specify the Laurent polynomial or parameter identification needed for this period comparison; those data and the precise meaning of singularity content are absent.
\item Equal quantum periods for class-\(\mathrm{TG}\) del Pezzo surfaces with the same qG-rigid singularities are asserted to imply qG-deformation equivalence. The required identification of period parameters and the meaning of having the same qG-rigid singularities are not specified. This clause of the original statement does not separately repeat the local qG-rigidity hypothesis imposed in the principal bijection.
\end{enumerate}
These companion assertions are source-reported formulations with unresolved conventions, not fully specified additional conjectures or proved consequences.
\subsection{Short English statement}
The main conjecture pairs mutation classes of Fano polygons with qG-deformation classes of locally qG-rigid class-\(\mathrm{TG}\) del Pezzo surfaces having cyclic quotient singularities, by generically deforming the associated toric surface. The original record also includes a mirror-map assertion and two period-detection assertions. Their missing conventions, and the unspecified ground field and deformation framework, require review. The supplied ``partially solved'' label is dated curation, not an independently verified mathematical conclusion.
\subsection{Sources}
\begin{itemize}
\item Supplied local UnsolvedMath registry record, permanent ID \(30002831\), title ``Mirror Classification of Orbifold Del Pezzo Surfaces'', exact source problem code \texttt{catalogue entry}, proposed year \(2015\). No individual proposer is specified.
\item Supplied report citation: \emph{Algebraic Geometry} (2015), Oberwolfach Reports, DOI \texttt{10.4171/owr/2015/15}, pp.~783--836. \url{https://doi.org/10.4171/owr/2015/15}. The record describes Oberwolfach Reports as documenting mathematical workshops and their open-problem discussions; no problem-specific page is supplied.
\item Akhtar--Coates--Galkin--Kasprzyk, supplied reference label ``original conjectures'': \url{https://kasprzyk.work/research/pdf/mirror_symmetry_del_pezzo.pdf}.
\item Cavey--Prince, supplied reference label ``verified cases'': \url{https://www.jstage.jst.go.jp/article/jmath/72/2/72_465/_pdf}.
\item Coates--Kasprzyk--Prince, supplied reference label ``later progress'': \url{https://pmc.ncbi.nlm.nih.gov/articles/PMC8526169/}.
\item Supplied literature assessment, checked \texttt{2026-08-21}: the record labels the problem ``partially solved'' and reports that the four mirror-symmetry conjectures are proved in smooth and several simplest singular cases, while the general mutation, period, and qG-deformation classification assertions remain conjectural. This is supplied dated curation derived from the OWR curation, not independently verified; its time-sensitive status and cited claims require independent checking.
\item The registry marks the clean statement \texttt{corrected\_verified} and reports that it was checked against the official Oberwolfach report for hypotheses, notation, and target. This is supplied verification metadata; no independent source consultation or mathematical verification is claimed here.
\end{itemize}
\end{document}
